\section{Thesis Structure}
\label{sec:introduction_thesis_structure}

The remainder of this thesis is structured as follows.

Chapter~\ref{chap:theoretical_background} introduces the theoretical and empirical background for the work. It discusses virtual embodiment, sense of agency, body ownership, self-location, visuomotor congruence, and related concepts such as presence and vection. It also reviews relevant work on VR exercise, sport simulation, and rowing-related systems. The chapter establishes the conceptual basis for treating movement mapping as a central design factor in VR rowing.

Chapter~\ref{chap:system_prototype} describes the VR rowing prototype developed for this thesis. It presents the hardware and software setup, the alignment and tracking approach, and the two representation modes. The ergometer-congruent mode and boat-centric mode are described in terms of their mapping strategies, avatar motion, and intended experiential differences. The chapter also discusses implementation constraints and design decisions that affect the comparison.

Chapter~\ref{chap:study_design} presents the study design used to evaluate the two representation modes. It describes the experimental conditions, participant procedure, measures, and analysis approach. Particular attention is given to how embodiment-related constructs, perceived rowing realism, comfort, and preference are measured. The chapter also explains how the comparison is structured to reduce confounds between the two modes.

Chapter~\ref{chap:results} reports the results of the evaluation. Quantitative and qualitative findings are presented with respect to the research questions, including measures related to agency, ownership, self-location, embodiment, perceived realism, comfort, and preference where applicable. The chapter focuses on reporting the observed outcomes without extensive interpretation.

Chapter~\ref{chap:discussion} interprets the findings in relation to the theoretical background and research questions. It discusses whether the results support the importance of visuomotor congruence, sport-realistic depiction, or a more nuanced relationship between the two. The chapter also considers limitations of the prototype and study design, as well as implications for future XR sport systems.

Chapter~\ref{chap:conclusion} concludes the thesis by summarizing the main findings and contributions. It revisits the central trade-off between ergometer-congruent and boat-centric representation and outlines directions for future work in VR rowing, embodied interaction, and XR sport design.
